\documentclass{article}
\usepackage[utf8]{inputenc}
\usepackage{xcolor}
\usepackage{hyperref}

\usepackage{rotating}

\usepackage{tabls}
\usepackage{booktabs}
\usepackage{amsmath}
\usepackage{amssymb}
\usepackage{amsthm}
\usepackage{amsfonts}
\usepackage{multicol}
\usepackage{enumitem}
\usepackage{microtype}
\usepackage{tikz}
\usepackage{pgfplots}
\usetikzlibrary{positioning}
\usepackage{multirow}

\usepackage{comment}

\usepackage{graphicx}
\usepackage{wrapfig}

\theoremstyle{plain}
\newtheorem{theorem}{Theorem}
\newtheorem*{theorem*}{Theorem}
\newtheorem{lemma}[theorem]{Lemma}
\newtheorem*{lemma*}{Lemma}
\newtheorem{proposition}[theorem]{Proposition}
\newtheorem{corollary}[theorem]{Corollary}

\theoremstyle{box}
\newtheorem{definition}[theorem]{Definition}
\newtheorem*{axiom*}{Axiom}
\newtheorem{dfn}[theorem]{Definition}
\newtheorem*{definition*}{Definition}
\newtheorem*{exercise*}{Exercise}
\newtheorem{observation}[theorem]{Observation}
\newtheorem{remark}[theorem]{Remark}
\newtheorem{example}[theorem]{Example}
\newtheorem{question}[theorem]{Question}
\newtheorem*{prep-problems}{Preparation Problems}

\newcommand{\pd}[3]{\frac{\partial^{#3} {#1}}{\partial {#2}^{#3}}}
\newcommand{\fpd}[2]{\frac{\partial {#1}}{\partial {#2}}}
\newcommand{\diff}[2]{\frac{d {#1}}{d {#2}}}


\begin{document}
In Project 2 we will be continuing our work with the LED data. To complete this project you will need to complete the following tasks insuring that you meet \textbf{ALL} the specifications. If \textbf{any} of the specification are not met you will receive feedback identifying the specification(s) that still require work in order to complete the project.

\begin{itemize}
\item Statistic  -- A number that summarizes information from data.
	\begin{itemize}
	\item Task 1a: Using the data for one light bulb, write down the expanded sums for each of the following statistics. Include the first three terms, ellipsis, and the last three terms. I  \textcolor{red}{(which light bulb can be randomly generated for each student)}
		\begin{enumerate}
		\item $\displaystyle \sum_{i=1}^{n} x_i$
		\item $\displaystyle \sum_{i=0}^{n} y_i$
		\item $\displaystyle \sum_{i=1}^{n} x_iy_i$
		\item $\displaystyle \sum_{i=1}^{n} x_i^2$
		\item $\displaystyle \sum_{i=1}^{n} y_i^2$
		\item $\displaystyle \frac{1}{n}\sum_{i=1}^{n} x_i$
		\item $\displaystyle \sum_{i=1}^{n} e^{-0.04x_i}$
		\item $\displaystyle \sum_{i=1}^{n} x_ie^{-2x_i}$
		\end{enumerate}
	\item Task 1b: Using the same data, calculate the value for each of the statistics listed in ``Task 1a".
	\end{itemize}
\item Optimize Likelihood Function
	\begin{itemize}
	\item Task 2a: Consider the function $f_1(t) = (100 - a_1) + a_1e^{-a_2t}$, write down the loglikelihood function.
	\item Task 2b: Write down the partial derivatives (for each parameter in the model) for the loglikelihood function.
	\item Task 2c: Set each partial derivative to zero and use R to numerically solve the resulting system of equations. Could this system of equations be solved analytically?
	\end{itemize}
\item Using the Fitted Model
	\begin{itemize}
	\item Task 3a: Create a plot with both the data (the scatterplot) and a fitted model from ``Task 2". Make sure you clearly indicate which parameter values you are using. Does the model visually fit the data?
	\item Task 3b: Use the model to predict when the luminescence of given light bulb will be less than 75\% of the original luminescence (luminescence at t = 0 hours).
	\end{itemize}
\item Assume a Different Functional Form for the Model
	\begin{itemize}
	\item Task 4a: Repeat ``Task 2" and ``Task 3" for the following models.
		\begin{itemize}
		\item $\displaystyle f_2(t) = (100 + a_1t)e^{-a_2t}$
		\item $\displaystyle f_0(t) = a_1 + a_2 t$
		\item $\displaystyle f_3(t) = 100 + a_2 t$
		\end{itemize}
	\item Task 4b: Compare and contrast the four fitted models you found.
	\item Task 4c: Explain your results in the context of the light bulbs.
	\end{itemize}
\item Conclusions
	\begin{itemize}
	\item Task 5: Answer the question, ``Do these light bulbs appear to meet the requirements for the L Prize?" Make sure to provide support for your answer using the models and your data. 
	\end{itemize}
\item Reflection
	\begin{itemize}
	\item Task 6: Reflect on your work for this project. In 1-2 paragraphs, identify/explain 2-3 key mathematical  ideas you learned (and would like to remember) and 1-3 soft skills you needed/improved/learned while working on the project.
	\end{itemize}
\end{itemize}


 
\end{document}